\sectionkicker{Theme synthesis}
\subsection*{横断比較 — bottleneckが移ると最適化層も移る}
\addcontentsline{toc}{subsection}{横断比較 — bottleneckが移ると最適化層も移る}
7月のserving更新は一つの高速化手法へ収束していない。runner、speculation、MoE kernel、schedulerという別々の層が同時に変わり、model architectureとhardwareの変化がruntime設計へ押し戻されている。
\medskip
\begin{center}
\small
\renewcommand{\arraystretch}{1.16}
\begin{tabularx}{\linewidth}{@{}p{0.20\linewidth}X@{}}
\toprule
観察軸 & 7月の一次資料・論文から読めること \\
\midrule
execution path & vLLMはdense model向けModel Runner V2をdefault execution pathへ移し、legacy PagedAttention implementationを削除した。後続releaseではattention、speculative decoding、KV offload、DeepSeek-V4関連optimizationも拡張された。 \cite{src-76c7b104c7df,src-b811cc97eff4} \\
speculative decoding & SGLangはSpec V2 pathをdefault化し、後続releaseでdraft confidenceに応じverification windowを変えるDSparkを導入した。 \cite{src-5a9041d41ad0,src-bf2a6d1b9b56} \\
MoE / hardware kernel & FlashInferはexpert-parallel MoE supportをdefault installへ入れ、MegaMoE、unified MoE API、Blackwell coverage、distributed servingを拡張した。 \cite{src-b80517d03022,src-020bbc9ce508} \\
research cost model & FlashDiffはdiffusion servingでregional execution/scheduling、EcoSpecはsparse MoE speculationでexpert activation/scattering costを扱う。いずれもauthor-evaluated researchであり、production superiorityではなく次の設計論点として読む。 \cite{src-50e29a09d75a,src-d6756f2b6aa6} \\
\bottomrule
\end{tabularx}
\renewcommand{\arraystretch}{1.0}
\end{center}
\normalsize
\begin{claimboundary}[この比較の境界]
この表は本号で既に検証・参照している一次資料と論文を横断比較の形へ再配置したものである。vendor / project / author が報告した性能値や優位性は独立再現済みの一般則へ変換せず、本文とSource Notesの帰属境界をそのまま維持する。
\end{claimboundary}
